\documentclass[aps,prb,onecolumn,nofootinbib,superscriptaddress]{revtex4-2}
\usepackage{amsmath,amssymb,amsthm,mathtools,bm}
\usepackage{booktabs}
\usepackage[colorlinks=true,linkcolor=blue,citecolor=blue,urlcolor=blue]{hyperref}
\usepackage{microtype}

\newcommand{\Ny}{N_y}
\newcommand{\Nx}{N_x}
\newcommand{\src}[1]{\texttt{\detokenize{#1}}}

\begin{document}

\title{Worker 09 Notes: Entanglement Scaling Evidence in the Adaptive Domain-Wall Ensemble}
\author{Worker 09}
\date{June 11, 2026}

\begin{abstract}
I summarize the repository evidence for logarithmic strip-entanglement scaling in the
large-$N_y$, $\Nx=20$ adaptive class-A domain-wall runs.  The inspected evidence supports
a compact but non-final claim: the trajectory-resolved perfect-correction ensemble produces
full-$x$ strip entropies with excellent log-chord form and slopes near the chiral target
$m=c/3=1/3$, with finite-size and finite-depth offsets.  The cleanest evidence is the
entropy sector; correlator exponents and chirality diagnostics are consistent with the same
domain-wall critical picture but remain less quantitatively settled.
\end{abstract}

\maketitle
\setcounter{tocdepth}{2}
\tableofcontents

\section{Source Scope}

I read \src{REPO_POLICY.md} before inspection.  Per the task constraints, I did not edit
files, execute notebooks, or load binary array products.  I excluded every path containing
\src{haining}, \src{haoyu}, or \src{erroneous}, case-insensitively, and did not read under
\src{repo_synthesis/main_pass/worker_notes}, \src{supervisor_notes}, or \src{boss}.

The main inspected sources were:
\begin{itemize}
\item \src{colab_large_entanglement_scaling_N20/README_COLAB.txt};
\item \src{colab_large_entanglement_scaling_N20/docs/entanglement_scaling_system_size.tex};
\item lightweight CSV/JSON summaries under
\src{colab_large_entanglement_scaling_N20/gpu_data/pure_state_entanglement_slope_vs_*};
\item \src{summary_of_results/summary_of_results.tex};
\item \src{summary_of_results/scripts/make_summary_figures.py}, as provenance for
\src{summary_of_results/figures/critical_scaling_summary.pdf}.
\end{itemize}

\section{Findings}

\subsection{Observable and CFT Target}

The measured observable is the von Neumann entropy of a full-$x$ strip of height $A_y$,
averaged over all starting positions $y_0$.  The fitting coordinate is the periodic
log-chord variable
\begin{equation}
X(A_y;\Ny)=
\log\left[
\frac{\Ny}{\pi}
\sin\left(\frac{\pi A_y}{\Ny}\right)
\right],
\end{equation}
and the fit is
\begin{equation}
\overline{S}(A_y)=m X(A_y;\Ny)+b.
\end{equation}
The local documentation identifies the domain-wall target as total central charge
$c=1$, hence
\begin{equation}
m_*=\frac{c}{3}=\frac{1}{3}.
\end{equation}
The fit window is $A_y\in[8,\Ny/2]$, intended to remove short-distance lattice points while
retaining the bulk of the log-chord curve.

\subsection{System-Size and Time Scaling}

The strongest finite-size/time-scaling note fixes $\Nx=20$, uses $\alpha_1=1$ in the
topological slab, $\alpha_2=30$ in the trivial exterior, $n_{\rm shell}=1$, and
\texttt{dw\_truncation=True}.  The documented time-scaling campaign uses
$\Ny=40,60,80,100,120$ with $T=\Ny/2$ Markov cycles and 10 samples per size.  The reported
slopes are
\[
0.3615,\ 0.3555,\ 0.3532,\ 0.3470,\ 0.3569
\]
for $\Ny=40,60,80,100,120$, respectively, with all $R^2\ge 0.9999$ except the quoted
$\Ny=120$ value still very high at $R^2=0.999907$.  The trend from $\Ny=40$ to $100$ is
monotone toward $1/3$, while $\Ny=120$ is a mild non-monotone outlier.

A separate system-size product at fixed cycle $50$ and 100 samples is described in the
Colab README for $\Ny\in[30,40,50,60,80,100,120]$, but the locally present
\src{pure_state_entanglement_slope_vs_system_size/runs} summaries visible in the inspected
tree include $\Ny=30,40,50$.  These give slopes $0.3557$, $0.3509$, and $0.3492$,
again above but close to $1/3$ with $R^2>0.99997$.

\subsection{Cycle Flow}

For the $\Nx=20,\Ny=40$ cycle sweep, the README states 100 samples and cycles $0,\ldots,100$.
The CSV shows the slope flowing rapidly from a nonuniversal early-time value to the
critical window.  For example, the first rows are extremely large at cycles $0$--$2$,
while cycles around $20$--$60$ sit near $0.35$ with excellent fit quality.  The same sweep
later drifts upward again near cycles $80$--$90$ before returning closer to the earlier
window by cycle $100$.  This makes ``cycle choice'' a real systematic, not a cosmetic
plotting detail.

\section{Interpretation}

The entropy evidence is best read as evidence for the shape and approximate coefficient of
a one-dimensional chiral critical edge sector embedded in the two-dimensional domain-wall
geometry.  The high $R^2$ values show that the full-$x$ strip entropy is not merely close to
the target at isolated points; it follows the log-chord curve very accurately over the
chosen window.  The remaining slope excess is plausibly finite-size and finite-depth
contamination, especially because the time-scaled $\Ny=40\to100$ sequence moves toward
$m=1/3$.

The broader synthesis note frames the relevant physical object as the trajectory-resolved
perfect-correction ensemble, not the post-selected branch or the mean covariance channel.
That matters because entropy and squared correlators are nonlinear observables.  In this
interpretation, the entanglement data provide the cleanest present evidence for the
quasi-one-dimensional chiral critical domain-wall ensemble, while correlator and modular
charge diagnostics are supporting but less mature.

The connection to chirality is indirect in the entropy fits alone: logarithmic scaling fixes
a critical sector and an effective central charge, not the propagation direction.  The
summary note adds modular-charge evidence that packets on opposite domain walls drift with
opposite signs, and it interprets \texttt{dw\_truncation=True} as a hard-wall projector
support limit of an inert trivial exterior rather than an independent source of chirality.
Those statements make the entanglement result part of a coherent chiral domain-wall story.

\section{Evidence Table}

\begin{table}[h]
\caption{Text and lightweight-summary evidence inspected for the entanglement scaling claim.}
\begin{ruledtabular}
\begin{tabular}{p{0.30\textwidth}p{0.14\textwidth}p{0.48\textwidth}}
Source & Lines & Evidence \\
\hline
\src{REPO_POLICY.md} & 50--58 & Required RevTeX shell and manuscript-note conventions. \\
\src{colab_large_entanglement_scaling_N20/README_COLAB.txt} & 7--14 & GPU runs use canonical \texttt{classA\_U1FGTN\_gpu.run\_markov\_circuit}; outputs are contour batches, entropy curves, fit rows, figures, metadata. \\
\src{colab_large_entanglement_scaling_N20/README_COLAB.txt} & 17--27 & Describes system-size and cycle notebooks, including $\Nx=20$, $n_{\rm shell}=1$, \texttt{dw\_truncation=True}, fit window $A_y=8..\Ny//2$. \\
\src{colab_large_entanglement_scaling_N20/README_COLAB.txt} & 29--32 & Domain-wall geometry for $\Nx=20$: wall locations $x=6,14$, topological slab $x=6..14$, trivial exterior regions. \\
\src{colab_large_entanglement_scaling_N20/docs/entanglement_scaling_system_size.tex} & 23--39 & Abstract-level claim: slopes $0.347$--$0.362$, target $1/3$, high $R^2$, finite-size/depth caveat. \\
\src{colab_large_entanglement_scaling_N20/docs/entanglement_scaling_system_size.tex} & 50--64 & Geometry, protocol, $\Nx=20$, $\Ny=40,60,80,100,120$, $T=\Ny/2$, 10 samples. \\
\src{colab_large_entanglement_scaling_N20/docs/entanglement_scaling_system_size.tex} & 68--90 & Defines full-$x$ strip entropy, log-chord scaling, target $m=1/3$, and fit window. \\
\src{colab_large_entanglement_scaling_N20/docs/entanglement_scaling_system_size.tex} & 103--118 & Table of time-scaled slopes and $R^2$ values for $\Ny=40$--$120$. \\
\src{colab_large_entanglement_scaling_N20/docs/entanglement_scaling_system_size.tex} & 137--162 & Interpretation: slopes within $4$--$9\%$, trend toward thermodynamic limit through $\Ny=100$, $\Ny=120$ non-monotonicity. \\
\src{colab_large_entanglement_scaling_N20/gpu_data/.../fit_rows.csv} & line 2 each & Fixed-cycle summaries: $\Ny=30,40,50$ slopes $0.3557,0.3509,0.3492$ at cycle 50. \\
\src{colab_large_entanglement_scaling_N20/gpu_data/..._time_scaling/runs/*/run_summary.json} & quoted fields & Time-scaled summaries record $T=\Ny/2$, slopes $0.3615,0.3555,0.3532,0.3470,0.3569$, target $1/3$, and $R^2$ values. \\
\src{colab_large_entanglement_scaling_N20/gpu_data/.../slope_vs_cycle.csv} & 1--25, 45--102 & Cycle sweep for $\Ny=40$: early nonuniversal decay into a near-$0.35$ slope window, later cycle dependence. \\
\src{summary_of_results/summary_of_results.tex} & 149--178 & States the target domain-wall critical theory: chiral complex fermion with $c=1$ and log-chord entropy scaling. \\
\src{summary_of_results/summary_of_results.tex} & 263--286 & Summary claim: entanglement fits approach $1/3$ more cleanly than squared-correlation exponents; correlators are not final exponent evidence. \\
\src{summary_of_results/scripts/make_summary_figures.py} & 373--427 & Provenance for \src{critical_scaling_summary.pdf}: reads slope-vs-cycle CSV, size fit rows, and correlation summary, then plots target lines at $1/3$ and exponent $2$. \\
\src{summary_of_results/summary_of_results.tex} & 316--348 & Interprets \texttt{dw\_truncation=True} as a hard-wall/infinite-mass exterior proxy, with caveats. \\
\src{summary_of_results/summary_of_results.tex} & 383--393 & Open issues: correlator extraction, ensemble ordering, modular-charge sizes/windows, and stable alpha-sweep tables. \\
\end{tabular}
\end{ruledtabular}
\end{table}

\section{Caveats}

The entropy evidence is strong for log-chord shape but not a completed thermodynamic
extrapolation.  The largest time-scaled point, $\Ny=120$, is non-monotone relative to
$\Ny=100$ and has visibly worse, though still excellent, $R^2$.  The note itself attributes
this to finite sample count and possible incomplete convergence at $T=60$.

There are two overlapping size campaigns in the available summaries.  The fixed-cycle
campaign is documented as broader, but the inspected local summaries under the main
system-size run directory expose only $\Ny=30,40,50$.  The time-scaled campaign supplies
the cleaner $\Ny=40,60,80,100,120$ progression but uses only 10 samples.

Cycle dependence is nontrivial.  The $\Ny=40$ slope is close to the target near cycles
$20$--$60$, but it is not strictly stationary through cycle $100$.  Any precision central
charge estimate should therefore specify both fit window and cycle-selection rule.

Finally, entropy scaling alone does not prove chirality.  It supports a critical sector
with effective $c\simeq1$ in the domain-wall geometry.  The chirality assignment relies on
the domain-wall construction, flattened-Hamiltonian calibration, and modular-charge drift
diagnostics summarized elsewhere.

\section{Confidence}

My confidence is high that the inspected repository evidence supports the qualitative
statement: the perfect-correction adaptive domain-wall ensemble shows clean logarithmic
log-chord entanglement scaling with slope near $1/3$.  My confidence is moderate for a
quantitative central-charge extraction because the accessible data still show finite-size,
finite-depth, and cycle-choice systematics.  My confidence is lower for using the current
correlator evidence as a precision universality check; the repository summary itself
treats those exponent fits as less settled than the entropy slopes.

\end{document}
